\section{Henry's law across the parameter space}
\label{s2:sec1}

\subsection{Henry's law holds everywhere; the region where it dominates does not}

Across all \nIsoNmat\ sampled materials, \(n > 1\) strictly for every material, so the cooperative term
is \(o(c)\) at the origin and the Henry constant is finite and positive everywhere.
But the Henry constant \(K_H = q_{\max} f_H b_H\), the slope of
equation~\eqref{eq:isotherm} at the origin, spans \nIsoKHmin--\nIsoKHmax\ \si{\mole\per\kilo\gram} per
\si{\mole\per\metre\cubed}, \nIsoKHdec\ decades, median \nIsoKHmed. The
concentration at which the cooperative term reaches 1\,\% of loading (the top of
the accessible Henry region) reaches \(10^{\nIsoHenryMax}\) of the feed at best
and a median of \(10^{\nIsoHenryMed}\), but lies below \(10^{-3}\) of the feed for
\nIsoHenryBelowMilli\ materials and below \(10^{-6}\) for \nIsoHenryBelowMicro. For
those the isotherm is effectively Sips-like at every concentration the column
visits. Since concentration is proportional to relative humidity at fixed
temperature, the near-linear uptake regime of a median material ends at
\(10^{\nIsoHenryMed}\) of the feed humidity. The design guarantees thermodynamic validity
everywhere, not an experimentally accessible linear regime.

\section{Where real water-harvesting frameworks sit}
\label{s2:sec2}

The sampled space is meant to
contain the adsorbents a water harvester would use, so we placed published ones in it
(Figure~\ref{fig:watermofs}), on the four axes that are properties of the framework
rather than of the pellet and the packing: step position, capacity, heat of
adsorption and the fraction of capacity taken below the step. The values come from
the primary literature, MOF-303 \cite{hanikel2019rapid,fathieh2018practical},
MOF-801 \cite{furukawa2014water}, CAU-10-H \cite{solovyeva2021cau,frohlich2016water},
aluminium fumarate \cite{hanikel2019rapid,laurenz2020novel}, MIL-160(Al)
\cite{solovyeva2021mil,aumond2026adsorption,lenzen2019metal} and
Co\(_2\)Cl\(_2\)(BTDD) \cite{rieth2017record}; each point spans the values its
sources report, which differ in sample, temperature and in how the step is defined,
and the below-step fractions are mostly read from published figures and are rough.
Of the \nWmofN\ frameworks, \nWmofInside\ lie inside the sampled ranges on every axis. The
exceptions are at the edges, and one is far outside: MOF-801's step, at a relative humidity of about
\nWmofEightStep, sits just below the sampled floor of \nWmofStepLo; one reported heat
for CAU-10-H, \nWmofCauHeat\,kJ\,mol\(^{-1}\), lies just beyond the sampled range; and
Co\(_2\)Cl\(_2\)(BTDD), at about \nWmofBtddQ\,mol\,kg\(^{-1}\), holds
\(\nWmofBtddRatio\times\) the largest capacity we sampled, \nWmofQmaxHi. The verdicts of this paper are therefore
statements about frameworks like MOF-303, aluminium fumarate and MIL-160, and should
not be carried to record-capacity materials without new data. The cooperativity
exponent \(n\), the steepness of the step and arguably the property a harvester
designer cares about most, is not placed: the sources report isotherms, not fitted
step exponents, and fitting equation~\eqref{eq:isotherm} to digitised curves would be
a new measurement, not a reading of the literature.

\begin{figure}[!tbp]
\centering
\includegraphics[width=\textwidth]{../figures/FigS2_water_mofs.pdf}
\caption{Published water-harvesting frameworks against the sampled parameter space
(shaded boxes). \textbf{A}~Step position and capacity. \textbf{B}~Fraction of capacity
taken below the step and heat of adsorption \(\Delta H\) (negative: adsorption is exothermic). Bars span the values reported across
sources; values, locations in each paper and confidence are recorded in
\texttt{research/water\_mof\_parameters.json}.}
\label{fig:watermofs}
\end{figure}

\section{The tuning envelope of rung L3}
\label{s2:sec3}

\paragraph{The tuning envelope is itself a result} The families optimise
\nLthreeLrDecades{} orders of magnitude apart: the Chebyshev family at \nLthreeChebyLr\ and the
perceptron at \nLthreeMlpLr. \nLthreeNfailAll{} configurations failed on every seed and
\nLthreeNfailSome{} more on some seeds; all are excluded and listed, never averaged in. Every family, the perceptron included, was tuned over the same
learning-rate grid; none was left at a default.

\section{Basis size on the legacy design}
\label{s2:sec4}

\begin{figure}[!tbp]
\centering
\includegraphics[width=\textwidth]{../figures/Fig4_operators.pdf}
\caption{Basis size on both designs. \textbf{A}~Legacy design (\nNclustLegacy\ held-out
materials): held-out error of each operator family against basis size \(p\), with the
basis floor; the Fourier operator has no basis size and is drawn as a level.
\textbf{B}~Full design (\nNclustFolds\ materials, five folds): DeepONet against \(p\) (one
faint point per fold, line through the means), the total error of POD with per-mode gradient boosting (basis floor included; the
map term alone is in the main text, Fig.~\ref{fig:basismap}), and the basis floor. The floor
falls by orders of magnitude; the total error does not move,
and DeepONet improves only from \(p = 8\) to \(p = 16\). \textbf{C}~Paired differences
against \(p = 8\), as a percentage of the error at \(p = 8\): on the legacy design
(grey) every interval spans zero; on the full design (black) every larger basis is
better. Per-pair intervals; the \(p = 128\) comparison also survives the
selection adjustment (A29). DeepOKAN, a DeepONet with Kolmogorov--Arnold layers; FNO,
Fourier neural operator.}
\label{s2:fig:operators}
\end{figure}


The legacy result below is retracted as a verdict (A29) and kept as a measurement. At
matched parameter counts and each family at its own best learning rate over
\nLfiveNrates\ rates spanning \nLfiveDecades\ decades, DeepONet looked flat in basis
size (Figure~\ref{s2:fig:operators}): \nLfiveFlatDiff\ from \(p = 8\) to \(p = 128\),
interval [\nLfiveFlatLo, \nLfiveFlatHi], which spans zero. Stated as what it excludes:
a sixteenfold increase in basis size could not have improved this arm by more than
\nLfiveFlatImproveMax\,\% of its error. The larger design found an improvement
well inside that bound, which is exactly why a null on \nNclustLegacy\ held-out
materials could not be the last word. Meanwhile its own POD
floor falls from \nLfiveFloorSmall\ to \nLfiveFloorLarge, a factor of
\(\nFloorDrop\). The best DeepONet found anywhere, at the smallest basis, sits
\(\nONetxFloor\times\) above the floor of the largest. The edge rule is the rule this
rung was most exposed on: at the two
smallest bases DeepONet's best learning rate first sat on the top of the grid, still
moving. The grid was extended one rate upward for those arms, the optimum stayed where
it was and is now interior, and every selected rate in this rung is bracketed, so
the Fourier operator of L5 is not compared with an under-tuned DeepONet. \emph{That DeepONet is as
accurate as an oracle handed the true coefficients of its first
\nLfiveOracleONet\ modes.}

\section{The scope of the L6 null}
\label{s2:sec5}

L6 predicted that the benefit of a separate kinetic object grows as the Damk\"ohler
number falls, and the main text finds no support for it. Kinetic control means Damk\"ohler well below unity, and this dataset never
goes there: its lowest per-material value is \nDaMatMin. The hypothesis predicted
that separating the kinetic object pays most where kinetics dominate, and we have
found no support for it over a range whose slow-kinetics end is merely \emph{less fast}, not
slow. A design that sampled \(\mathrm{Da} < 1\) could in principle recover the
interaction. We did not sample it, and the reason is structural rather than an
oversight: the kinetic coefficient is not a free parameter in this dataset but is
derived from particle size through the Glueckauf relation, so
\(\mathrm{Da} \propto d_p^{-2}\) and reaching \(\mathrm{Da} = 1\) from our slowest
bed would take a pellet of about \nDaPelletForUnity\,\si{\milli\metre}. That is
\nDaPelletsAcross\ pellet diameters along the \nDataBedLenCm\,\si{\centi\metre} column, where a
one-dimensional plug-flow model is no longer the right description. This is the same
objection Section~\ref{sec:discussion} raises against our own anchor bed. Decoupling
\(k_{\mathrm{LDF}}\) from \(d_p\) to reach that regime would restore precisely the
physical inconsistency this dataset was built to remove. The null therefore
holds across the range a physically consistent bed of this kind can occupy, and is
not a statement about the kinetically-controlled limit.

\section{The power of the L6 test}
\label{s2:sec6}

\paragraph{Why the test has power now, and its limit} The between-material standard
deviation of the arm difference fell to \nLsixBetweenSD, or
\nLsixBetweenOverBase\,\% of the base error, from \nLsixLegacyBetweenOverBase\,\% previously, a \(\nLsixNoiseDrop\times\) noise reduction from Sobol sampling and physically consistent
kinetics. With \nNclustFolds\ clusters instead of \nNclustLegacy, that is what turned an unanswerable
question into an answerable one. Power is \nLsixPowerFour\,\% at a 4\,\% effect and
\nLsixPowerTwo\,\% at 2\,\%, with size at the null \nLsixSizeNull\,\%; the minimum
detectable effect at 80\,\% power is \nLsixMDE\,\%. The observed effect sits just
under it, so the point estimate may be optimistic and the lower bound is the
conservative reading.

\section{The mode count and the n-width}
\label{s2:sec7}

\paragraph{The mode count is scoped to a sample size} An
earlier version of the paper stated the mode count as a property of the map. It
is not: it is what a given number of training materials buys. Measured against
material count (Figure~\ref{fig:basismap}), modes clearing \(R^2 > 0.2\) grow from \nLCmodesSmall\ to
\nLCmodesLarge\ over 12 to 192 materials and are still climbing. Three counts of
modes clearing \(R^2 > 0.2\) appear in this paper, and they are the same statistic
under different conditions: \nLfiveCountModes\ on the legacy design, and
\nLCmodesLarge\ and \nLfiveVtwoModes\ on the larger one at the same number of training
materials, the first for the learning-curve learner on its own basis and the second for the
decomposition's regressor on the largest basis. They agree on what matters here:
a learner predicts a modest set of mostly leading modes and little beyond it. The retraction is
A21, and it withdrew the replacement text for an earlier retraction: the same
inference, made twice.

\paragraph{The \(n\)-width is real and was not what binds} The residual
\emph{energy} of this snapshot family decays algebraically, as
\(k^{\nNwidthExpC}\) in the number of modes \(k\), with \(R^2 = \nNwidthRtwoC\) over
\(k = 8\)--64. The Kolmogorov \(n\)-width, whose \(n\) is this mode count, is the
residual \emph{norm}, the square root of that, so it decays as \(k^{\nNwidthExpErr}\)
and halving the error requires
\(\nNwidthHalving\times\) the modes. We separate the two because an
earlier draft quoted the energy exponent and called it the \(n\)-width, next to a
halving factor computed correctly from the norm: two numbers in one sentence
that disagreed by a factor of two in the exponent. Any method whose output is a linear
combination of a fixed basis is bounded below by it. Our frozen protocol called this
the central result of the ladder and predicted that DeepONet would track it.
The prediction was wrong: the bound is real, and it is inactive. Withdrawing it is
A17. For transport-dominated problems the slow decay is established: Ohlberger
and Rave \cite{ohlberger2016reduced} prove a \emph{lower} bound of \(\tfrac12 k^{-1/2}\) for linear advection
with jump discontinuities, which is a bound and not a rate. Our measured decay is faster than that bound, and the reason is the
hypothesis: their fields carry jump discontinuities, ours are smoothed by axial
dispersion and by finite-rate kinetics, so the worst case their theorem
constructs does not occur here. That is also why the bound can be simultaneously
real and inactive; this paragraph has to hold both statements together.

\section{Monotonicity of the predicted fronts}
\label{s2:sec8}

The two-front frame of Section~\ref{s2:sec:warp} predicts two front trajectories per
sample. The obvious first attack on its result is that the front predictor does not know the constraints the true trajectories satisfy exactly.
It does violate them, in \nWarpNonmonoLo\,\% of samples at the lower landmark, while the true trajectories
violate them in none. But projecting onto the constraint set changes \(R^2\) by
at most \nWarpMonoMaxDRtwo{} and makes the reconstruction significantly
\emph{worse}: \nMonoRawNrmse\ against \nMonoProjNrmse, a difference of
\nMonoDiff\ with interval [\nMonoLo, \nMonoHi] (tighter than Table~\ref{s2:tab:warp}'s
because the projected arm is a deterministic transformation of the raw arm, so the
between-material variance cancels). The violations are frequent but negligible; the
remaining headroom is the front predictor's raw accuracy.

\section{The dispersion closure, and the limits of the isotherm and the parameter space}
\label{s2:sec9}

Our datasets use
\(D_L = 0.7 D_m + 0.5 d_p u\), with \(u = v/\varepsilon_t\) the interstitial velocity
and \(D_m\) the molecular diffusivity of water in air, a form due to Wakao
and Funazkri \cite{wakao1978effect} that Perry's handbook \cite{levan2019adsorption}
presents as a \emph{lower bound} for \emph{nonporous} particles; for porous, strongly
adsorbing particles such as ours the effective axial dispersion is larger. The form is
also the high-particle-P\'eclet \emph{limit} of Edwards and Richardson's correlation
\cite{edwards1968gas},
\(D_L = 0.73 D_m + 0.5\,d_p u\left[1 + 9.7\,D_m/(d_p u)\right]^{-1}\),
whose mechanical term carries the extra bracketed factor, so we measured whether the
limit applies.
The particle P\'eclet number \(d_p u / D_m\) runs
\nPePartMin--\nPePartMax\ across the dataset, median \nPePartMed, spanning
\nPePartDecades{} decades. Over that range the two published forms disagree by a
median factor of \nDispRatioMed, exceeding \(1.1\times\) for \nDispFracAboveTenPct\,\%
of samples and \(1.5\times\) for \nDispFracAboveHalf\,\%, with a worst case of
\nDispRatioMax. The disagreement is bounded and modest for the dataset, but it is
not modest for the anchor.

\begin{figure}[!tbp]
\centering
\includegraphics[width=\textwidth]{../figures/Fig7b_anchor_posthoc.pdf}
\caption{Post-hoc dispersion sensitivity on the anchor bed, labelled as such and
reported as a limitation rather than as a result. The pre-registered primary uses the
high-P\'eclet limit of the correlation; the second curve is the \emph{same}
correlation without that truncation, at the same particle size, velocity and
molecular diffusivity; the third is the authors' own Bruggeman closure for the
effective dispersion \(D_e\); the fourth repeats that closure at the smallest LDF
constant \(k\) (in s\(^{-1}\)) that still keeps the step. Pe is the column P\'eclet
number. The measured 95\,\% arrival time is bracketed by the untruncated correlation and
the authors' Bruggeman closure.}
\label{fig:anchorposthoc}
\end{figure}

The anchor bed sits at a particle P\'eclet number of \nAnchorPePart, and the
disagreement between the two forms is maximised at \nDispWorstPe, where it reaches
\nDispWorstRatio\(\times\). In other words the one bed in this paper that is compared
against an experiment is also the one place where the truncation matters most: there
our closure disperses \nAnchorDispRatio\(\times\) more than the correlation it is the
limit of, taking the column P\'eclet number from \nAnchorPeColER\ to \nAnchorPeCol.
Re-solving that bed under the untruncated form (a labelled post-hoc sensitivity,
not a fit, and not the pre-registered result) moves the 95\,\% arrival time from
\nAnchorTninetyfiveModel\ to \nAnchorTninetyfiveER\ minutes against
\nAnchorTninetyfiveExp\ measured, the plateau from \nAnchorPlateauModel\ to
\nAnchorPlateauER\ against \nAnchorPlateauExp, and the error against the digitised
effluent from \nAnchorNrmse\ to \nAnchorNrmseER\ (Figure~\ref{fig:anchorposthoc}).
The authors' own Bruggeman closure goes further, to \nAnchorNrmseBrug, but overshoots
in the other direction at \nAnchorTninetyfiveBrug\ minutes; the measured value is
bracketed by the two closures. A substantial part of the anchor's second named
discrepancy is therefore one term evaluated at its limit: the untruncated form
removes \nAnchorTcutPct\,\% of the model's excess 95\,\% arrival time and
\nAnchorNrmseCutPct\,\% of its error against the digitised effluent (the error is
dominated by the first-wave timing, which no dispersion closure touches). None of this
affects a surrogate result, because every arm sees the same reference; it bounds how far
the reference may be read as physical.

\paragraph{The isotherm's low-humidity branch is too weak} The Henry fraction was
fitted to the step position and the capacity, never to dynamics; the anchor comparison
shows the model holds too little water below the step.

\paragraph{Part of the heat range is occupied by no placed framework} The sampled
heat of adsorption runs from \nDcDHLo\ to \nDcDHHi\,kJ\,mol\(^{-1}\); the
water-harvesting frameworks we placed (Figure~\ref{fig:watermofs}) report \nWmofHeatMost\ to
\nWmofHeatLeast. Near its least exothermic edge the sampled range gives an isosteric
heat (\(-\Delta H + RT\)) below water's latent heat of vaporisation at room
temperature, so that water is bound more weakly than in its own liquid. That is the
regime of hydrophobic adsorbents rather than of the frameworks this paper is
motivated by. The range was pre-registered and is kept; verdicts averaged over it
weight that region as heavily as the occupied one.

\paragraph{Not every sample reaches the step} A harvester is operated above its
step by design; the sampled space is not. \nChemBelowStepPct\,\% of the accepted
samples of the larger dataset (\nChemBelowStepLegPct\,\% of the legacy one) are fed
below their material's step humidity, so the feed never reaches pore filling and the
breakthrough is the primary-site front alone. These samples are in every verdict; a
surrogate built only for harvesting conditions would not need them.

\paragraph{The step does not move between feed temperatures} The cluster affinity is
set per run so that the step sits at the material's step humidity at the feed
temperature (Section~\ref{sec:model}). Within a run the step still moves with the
local bed temperature through the heat of adsorption, but across runs a material's
step humidity at the feed temperature is the same by construction. In real MOF-303
it moves with temperature \cite{hastings2024high}. The
datasets therefore do not test whether a surrogate transfers
the temperature dependence of the step.

\paragraph{Three parameters are not measurements} Pellet density and bed porosity
are packing values, not MOF-303 measurements \cite{hanikel2021evolution}. No citable measured heat capacity for
MOF-303 was found, so it is reported as a sensitivity band rather than as a cited
constant.

\paragraph{No covalent organic framework case} We found no published breakthrough
measurement on a covalent organic framework (COF) suitable for this comparison, so no
claim is made for that class.

\section{When each pre-registration was fixed}
\label{s2:prereg}

Each pre-registration is a file in the public repository, fixed by the commit that added
it. Table~\ref{tab:prereg} gives, for each, that commit and the first commit of the results
it governs, by author date. The repository's history was rewritten once, on 5 October 2026,
before it was made public, to remove third-party files; the rewrite kept every author date
and changed every commit identifier, and \texttt{COMMIT\_MAP.tsv} maps each old identifier
to its new one. The ordering below is therefore attested by the author's own history, not by
an external registry. The original ladder protocol is the weakest case: it was written before
its runs but committed only when the repository was created, minutes before the first
results files, so for the rungs it governs (L3, the first L5 sweep and the coordinate
change, all on the exploratory 12-material design) the order rests on the author's word. The
two additions made for the Digital Discovery revision were pushed to the public repository
before they ran, so the push is an external timestamp.

\begin{table}[!htbp]
\centering
\caption{Pre-registrations and the first results they govern (author dates, UTC+3).}
\label{tab:prereg}
\small
\begin{tabular}{@{}>{\raggedright\arraybackslash}p{4.6cm}>{\raggedright\arraybackslash}p{3.9cm}
                >{\raggedright\arraybackslash}p{2.7cm}>{\raggedright\arraybackslash}p{3.0cm}@{}}
\toprule
pre-registration & governs & fixed & first results \\
\midrule
\texttt{03\_LADDER\_PROTOCOL.md} & L3, first L5 sweep, coordinate change & 30 Aug 2026, 17:50 & 30 Aug 2026, 17:57 \\
\texttt{PREREG\_L6\_v2.md} & L6 & 1 Sep 2026 & 2 Sep 2026 \\
\texttt{PREREG\_L1L2L7\_v2.md} & L1, L2, L7 & 2 Sep 2026 & 3 Sep 2026 (L1, L7), 5 Sep (L2) \\
\texttt{PREREG\_L4b\_v2.md} & L4 & 2 Sep 2026 & 24 Sep 2026; corrected re-run 11 Oct \\
\texttt{PREREG\_LASSITTER.md} & the MOF-303 anchor & 3 Sep 2026, 09:56 & 3 Sep 2026, 10:03 \\
\texttt{PREREG\_L5\_v2.md} & L5 on 240 materials & 28 Sep 2026 & 4 Oct 2026 \\
\texttt{PREREG\_BASELINES.md} & Gaussian process, coarse solve, fitted closed forms & 4 Oct 2026 & 11 Oct 2026 \\
\texttt{PREREG\_DD\_ADDITIONS.md} & screening, named frameworks & 25 Oct 2026 (pushed) & 25 Oct 2026, after the push \\
\bottomrule
\end{tabular}
\end{table}

